\documentclass[11pt]{article}
\usepackage[utf8]{inputenc}
\usepackage[T1]{fontenc}
\usepackage{graphicx}
\usepackage{xcolor}
\usepackage{amsmath, amssymb}
\usepackage{url}
\usepackage{natbib}
\setlength{\parindent}{0em}
\usepackage[letterpaper, total={6.5in, 9in}]{geometry}

\newcommand*{\mybox}[2]{%
  \par\medskip\noindent
  \colorbox{#1!30}{\parbox{.98\linewidth}{#2}}%
  \par\medskip
}

\setlength{\parskip}{0.5em}

\usepackage{setspace}
\onehalfspacing

\usepackage{xr-hyper}   
\makeatletter
\newcommand{\myexternaldocument}[2][]{%
  \externaldocument[#1]{#2}%
  \IfFileExists{#2.aux}{}{%
    \typeout{No file #2.aux.}%  
  }%
}
\IfFileExists{newcommands.tex}{% --- LaTeX Packages ---
% --- LaTeX Packages ---
\usepackage{amsmath}
\usepackage{amsthm}
\usepackage{bm}            % For bold math symbols
\usepackage{graphicx}      % For including images
\usepackage{xcolor}        % Required for \highlight
\usepackage{xspace}        % Required for commands like \nref
\usepackage{tikz}
%\usetikzlibrary{arrows.meta}
%\usepackage{mathtools}
%\usepackage{filecontents}
\usepackage{booktabs}
%\usepackage{placeins}
\usepackage{hyperref}
        \hypersetup{
          colorlinks=true,
          linkcolor=blue,
          urlcolor=blue,
          citecolor=blue
        }
\usepackage{etoolbox}

\setlength{\bibsep}{5pt plus 5pt} % space BETWEEN items
% For tighter lines WITHIN each item:
\usepackage{setspace}
\AtBeginEnvironment{thebibliography}{\setstretch{1.05}}

% --- Custom Commands ---

% Acronyms and Custom Text
\newcommand{\CVIC}{\text{CVIC}}
\newcommand{\deviance}{\text{dev}}
\newcommand{\dpois}{\text{dpois}}
\newcommand{\ic}{\text{ic}}
\newcommand{\nmsp}{\text{\scriptsize NMSP}}
\newcommand{\nrsp}{\text{\scriptsize NRSP}}
\newcommand{\pmin}{p_{\text{\scriptsize min}}}
\newcommand{\pp}{\text{pp}}
\newcommand{\nref}{[\textbf{references}]}
\newcommand{\SP}{survival probability\xspace}
\newcommand{\SPs}{survival probabilities\xspace}
\newcommand{\svf}{S_{i}(\cdot)}
\newcommand{\unif}[0]{\text{Unif}}
\newcommand{\rsp}[0]{\text{RSP}}
\newcommand{\msp}[0]{\text{MSP}}
\newcommand{\PPP}[0]{\text{PPP}}
\newcommand{\PPPs}[0]{\text{PPPs}}
\newcommand{\PPC}[0]{\text{PPC}}
\newcommand{\bU}[0]{\bm{U}}
\newcommand{\pv}[0]{\mathcal{P}}

% Math Symbols (Bold)
\newcommand{\btheta}[0]{\bm{\theta}}
\newcommand{\bTheta}[0]{\bm{\Theta}}
\newcommand{\bx}[0]{\bm{x}}
\newcommand{\by}[0]{\bm{y}}
\newcommand{\bY}[0]{\bm{Y}}
\newcommand{\bbeta}{\bm{\beta}}
\newcommand{\bBeta}{\bm{\Beta}}
\newcommand{\blambda}{\bm{\lambda}}
\newcommand{\bmu}{\bm{\mu}}
\newcommand{\bphi}{\bm{\phi}}
\newcommand{\bp}{\bm{p}}
\newcommand{\bsgm}{\bm{\sigma}}
\newcommand{\hu}[0]{\pi^o_{i}}
\newcommand{\mhu}[0]{\pi^o}
% General Superscripts and Subscripts
\newcommand{\obs}[0]{^{\text{obs}}}
\newcommand{\sobs}[0]{^{\text{obs}}}
\newcommand{\post}[0]{^{\text{Post}}}
\newcommand{\cv}[0]{^{\text{CV}}}
\newcommand{\iscv}[0]{^{\text{ISCV}}}
\newcommand{\supt}[0]{^{(t)}}
\newcommand{\mci}{^{(t)}}        % Kept as it has a different name from \supt
\newcommand{\mcs}{^{(s)}}
\newcommand{\rep}{^{\text{rep}}}
\newcommand{\ppc}{_{\text{ppc}}}
\newcommand{\noi}{_{-i}}
\newcommand{\wi}{_{1:n}}
\newcommand{\IS}{^{\text{\scriptsize IS}}}

% Compound Superscripts for Z-residuals
\newcommand{\rspp}[0]{^{\text{Post-RSP}}}
\newcommand{\mspp}[0]{^{\text{Post-MSP}}}
\newcommand{\rspis}[0]{^{\text{ISCV-RSP}}}
\newcommand{\mspis}[0]{^{\text{ISCV-MSP}}}
\newcommand{\rsppost}{^{\text{Post-RSP }}}  % From first block
\newcommand{\rspiscv}{^{\text{ISCV-RSP}}}    % From first block
\newcommand{\E}{\mathbb{E}}
\newcommand{\Epost}{\mathbb{E}_{\btheta\,|\,\by}}
\newcommand{\Ecv}{\mathbb{E}_{\btheta\,|\,\by_{-i}}}
\newcommand{\highlight}[1]{\textcolor{blue}{#1}}
\newcommand{\pvalue}[0]{$p$-value\xspace}
\newcommand{\pvalues}[0]{$p$-values\xspace}
\newcommand{\norm}[1]{\left\lVert#1\right\rVert}

\newcommand{\textred}[1]{\textcolor{orange}{#1}}


% ===================================================================
%                 THEOREM AND PROOF DEFINITIONS
% ===================================================================



%old theorem style for ba

% % --- 1. Define Custom Theorem Styles ---
% % A style for theorems where the optional note is bold
% \newtheoremstyle{boldnote}
%   {}		    % Space above
%   {}		    % Space below
%   {\itshape}	% Body font
%   {}		    % Indent amount
%   {\bfseries}	% Theorem head font
%   {.}		    % Punctuation after head
%   {.5em}	    % Space after head
%   {\thmname{#1}\thmnumber{ #2}\thmnote{ [\textbf{#3}]}}


% % --- 2. Define Theorem-Like Environments ---
% % Set the default style for all environments defined from this point
% \theoremstyle{plain}

% % The main theorem environment, numbered per section. All others will follow its counter.
% \newtheorem{theorem}{Theorem}[section] 

% % Environments that share the theorem counter
% \newtheorem{lemma}[theorem]{Lemma}
% \newtheorem{corollary}[theorem]{Corollary}

% % Switch to the custom style for any environments you want to use it
% % \theoremstyle{boldnote}
% % \newtheorem{specialtheorem}[theorem]{Special Theorem} % Example


% % --- 3. Define Proof and Step Environments ---
% % For custom "Step" counter within proofs
% \newtheorem{proofpart}{Step}

% % Change the name "Proof" to "Proof" (bold)
% \renewcommand{\proofname}{\textbf{Proof}}

% % Automatically reset the 'Step' counter at the beginning of every proof
% \AtBeginEnvironment{proof}{\setcounter{proofpart}{0}}
% % ===================================================================
}{}

\makeatother


\myexternaldocument[S-]{supplement}
\myexternaldocument[M-]{JASA-template}

\makeatother



\begin{document}

\begin{center}
    {\LARGE \textbf{Response to Reviewer C}} \\[1em]
    {\large \today}
\end{center}
\vspace{2em} % Space between title and your text
\normalsize



We sincerely thank you for your time, valuable feedback, and general appreciation of the paper. Guided by the collective insights from you and other referees, we have substantially revised the manuscript. We have responded to your comments item by item, and your constructive suggestions have greatly improved both the methodology and the presentation. In this detailed response, we quote several updated passages for your reference (please note that some of these quoted excerpts may have been further refined in the final manuscript to ensure perfect flow and accuracy). In the revised manuscript, revised text is highlighted in \textred{orange}; for substantially rewritten sections, only the \textred{section title} is highlighted rather than the entire text. 

\mybox{gray}{\textbf{Summary.} The paper presents an approach to Bayesian model checking based on the notion of $Z$-residuals, which are defined for a given observation $i$ as
\[
z_i^{\mathrm{RSP}}(y_i \mid \theta) = -\Phi^{-1}\!\bigl(\mathrm{RSP}_i(y_i \mid \theta)\bigr),
\]
where $\Phi(\cdot)$ is the cdf of a standard normal and $\mathrm{RSP}_i(y_i \mid \theta)$ is the randomized survival probability, defined as
\[
\mathrm{RSP}_i(y_i \mid \theta) = S_i(y_i \mid \theta) + U_i p_i(y_i \mid \theta),
\]
where $U_i \sim \mathrm{Unif}(0,1)$, $S_i(y_i \mid \theta)$ is the survival function of $Y$, and $p_i(y_i \mid \theta)$ is its probability mass function (equal to zero if the random variable is continuous). Specifically, the authors show that if the model is correctly specified, then $\mathrm{RSP}_i(y_i \mid \theta) \sim \mathrm{Unif}(0,1)$. In turn, this makes $z_i^{\mathrm{RSP}}(y_i \mid \theta)$ normally distributed. In particular, the authors use this property to construct posterior $Z$-residuals, obtained by averaging $\mathrm{RSP}_i(y_i \mid \theta^{(t)})$ for each posterior sample $\theta^{(t)}$. They also introduce a leave-one-out cross-validated version paired with importance sampling, to further correct the problematic double use of the data (which undermines calibration of the $p$-values). Their approach is useful for testing goodness-of-fit, e.g., via a $\chi^2$ test, as well as a tool to capture potential misspecifications in the functional form of a covariate. The authors showcase their approach in a variety of simulated settings and in a real dataset of dolphin sighting. \\[5pt]
\textbf{General comment.} This is an interesting and potentially helpful contribution. While the notion of RSP is equivalent to the probability integral transform, which has been used for model checking in the Bayesian literature, what the author presents is a general strategy for checking both the correct specification of the likelihood and also which covariates might be misspecified. Moreover, the theoretical statements and proof appear correct. However, I have the following comments and questions.
}

\paragraph{Response:}

We sincerely thank the reviewer for this excellent and highly accurate summary of our manuscript. The reviewer's understanding of the proposed Z-residual framework, the role of randomized survival probabilities (RSPs), and the posterior and ISCV implementations is perfectly correct.

We are particularly grateful for the reviewer's recognition of the core innovation of this work. As the reviewer acutely pointed out, while the mathematical foundation of RSPs is equivalent to the probability integral transform in the context of continuous response, our primary contribution is elevating this concept into a comprehensive, general strategy for devising versatile normally-distributed residuals that enable the diagnosis of  overall likelihood specification, localized covariate misspecifications, and many other model deficiencies of the likelihood. We also greatly appreciate the positive assessment of our theoretical statements and proofs.

Inspired by this summary, we have substantially revised the Section of Introduction to articulate these contributions in a clearer way: it now more clearly distinguishes the PIT foundation shared with the existing literature from the new elements of our framework---the summarize-first construction of a single residual set, its asymptotic calibration, and the resulting suite of targeted graphical and numerical diagnostics. 


\medskip

\mybox{gray}{\textbf{Question 1.} The author presents differences between the RSP, MSP, and PPC. However, they should also compare their approach against the more recent holdout predictive checks (HPC) of Moran et al.\ (2023), which also avoid using the data twice and provide calibrated $p$-values, at the cost, however, of lower power. Since the proposed approach does not suffer from this loss of power, it would be great to see examples where your method works better (or comparatively).}

\paragraph{Response:}
Thank you for this excellent advice. We have included holdout predictive checks (HPC) in all of our empirical studies, and we believe that adding this benchmark has greatly strengthened the paper. 

Interestingly, our empirical studies additionally revealed that HPC is not only less powerful but can also exhibit slightly worse-calibrated Type I errors in small-sample settings. While initially surprising, this behavior is reasonable upon reflection: the finite-sample posterior of $\theta$ (trained on the observed data) is not identical to the true parameter governing the isolated test data. Because the posterior predictive distribution integrates over this finite-sample posterior uncertainty, it is structurally different from the true data-generating distribution. Consequently, evaluating holdout data against this approximate predictive distribution can distort the exact uniformity of the resulting $p$-values in small samples. 

We have revised Sections 3 and 4, as well as Supplementary Section S1, to include the HPC comparisons and to discuss these finite-sample power and calibration trade-offs explicitly.


\medskip

\mybox{gray}{\textbf{Question 2.} It would be interesting to discuss your $Z$-residuals with the recent U-values posed by Covington and Miller (2025). Their approach introduces uniform distributions for each of the model parameters, while yours adds a $z$ variable for each observation. Under which settings should your approach be preferred?}

\paragraph{Response:}
Thank you for this helpful suggestion. Discrete response is essentially an incomplete observation of an underlying continuous random variable. The U-value approach of Covington and Miller (2025) is closely related in spirit to ours: both transform random elements onto a uniform reference scale for recovering the underlying continuous random variables. Indeed, for a single posterior draw $\btheta^{(t)}$, their observation-level uniform variables are distributionally identical to $1 - \rsp_i(y_i \mid \btheta^{(t)})$. The two frameworks differ in two key respects. First, in posterior aggregation: UPC follows a ``test-first, then-summarize'' strategy, combining MCMC-wise $p$-values via the Cauchy combination rule \citep{liu2020cauchy}, whose approximation error may not fully vanish asymptotically, whereas we summarize RSPs over the posterior or leave-one-out predictive distribution first, yielding a single residual set that is asymptotically i.i.d.\ $N(0,1)$. Second, in diagnostic target: U-values can be attached to any random element of the generative specification---parameters, priors, or latent variables---making them well suited to locating which internal component of the model is problematic. Z-residuals, in contrast, are observation-level predictive residuals on a standard normal scale, supporting regression-style diagnostics such as QQ plots, residual-versus-covariate plots, grouped boxplots, and functional-form tests.

Accordingly, our approach is preferable when the goal is to assess the fitted predictive distribution at the observation level---for example, tail behavior, nonlinear or omitted covariate effects, or zero versus nonzero patterns in count models---while U-values are more appropriate for questions about parameter-, prior-, or latent-variable components. We view the two approaches as complementary rather than competing, and we have added a detailed discussion of Covington and Miller (2025) to Section 5 and a discussion of how they can be combined to achieve the diagnosis of internal constructs as are demanded in applied Bayesian modelling.

\medskip

\mybox{gray}{\textbf{Question 3.} The authors indicate that the $Z$-residuals may detect potential misspecification in the functional form of some of the covariates (e.g., the simulation in Section 3.3). By looking at Figure 4(c), it looks like the trend between the $z$-residuals and $x_1$ in the model that assumes $x_1$ as linear has approximately a log-shape, which is the true functional form.
}

\paragraph{Response:}
Thank you for this acute observation, which is consistent with our experience. When the response family is correctly specified, the scatterplot of Z-residuals against a covariate under a misspecified linear model often reflects the shape of the true functional form---in this case, the logarithmic shape visible in Figure~4(c). Intuitively, the Z-residuals capture the discrepancy between the true covariate effect and its best linear approximation, so the residual trend traces the unmodeled portion of the true curve. This property makes the residual-versus-covariate plot not only a detector of functional-form misspecification but also a practical guide for choosing a corrected specification, as illustrated in our dolphin-sighting application, where the residual pattern guided the revision to the nonlinear HNB model, resulting in a more predictive model.
\mybox{gray}{
\textbf{3(a)} The HNB family here is correctly specified. What would happen if one were to model the data as an NB model, under both the correct functional form for $x_1$ (logarithmic) and the linear one? It would be important to clarify if the potential functional form between a given $x$ and the $z$-residuals is impacted by the misspecification of the likelihood. A hint on what happens is suggested in Table 2 in the application (NB vs HNB in the nonlinear part), but since it is real data, it is hard to grasp the true relationship.}

\paragraph{Response:}
Thank you for this very thought-provoking question. Following the
reviewer's suggestion, we have added simulation studies in the new
Supplementary Section~S1.2, using the same HNB-generated data as in
Section~3.3 and fitting two additional NB models: NB-log (misspecified
response family) and NB-linear (both misspecified). The Z-residual
scatterplots against $X_1$ and the diagnostic test results for a
representative dataset are presented in Section~S1.2; the $p$-values are
reproduced in the table below.

\begin{table}[!htbp]
\centering
\caption{Z-residual diagnostic $p$-values for the response-family sensitivity analysis using the same HNB-generated representative dataset as in Section~\ref{M-sec:sim-nl} of the main text. The HNB-log and HNB-linear rows correspond to the models considered in the main text, while the NB-log and NB-linear rows are added here to assess the effect of response-family misspecification. Because the response is discrete and RSP Z-residuals depend on auxiliary randomization, we report median upper-bound $p$-values~(Eq.~\eqref{M-eqn:uppvalue-median}) for the posterior and ISCV RSP Z-residual tests. The AOV and Bartlett (BL) tests are applied against $X_1$. The logarithmic form is marked ``Wrong'' for the NB family because the $X_1$-dependent hurdle probability is improperly absorbed into the NB mean structure (see text).}
\label{tab:sens_family_function}
\scriptsize
\resizebox{\textwidth}{!}{%
\begin{tabular}{lllrrrrrrrr}
\toprule
 & & & \multicolumn{4}{c}{Posterior RSP Z-residuals} & \multicolumn{4}{c}{ISCV RSP Z-residuals} \\
\cmidrule(lr){4-7} \cmidrule(lr){8-11}
Model & Family & Functional form 
& SW & AOV-$X_1$ & BL-$X_1$ & $\chi^2$
& SW & AOV-$X_1$ & BL-$X_1$ & $\chi^2$ \\
\midrule
HNB-log    
& HNB (True) & log (True)   
& 0.904 & 0.630 & 0.120 & 0.884 
& 0.595 & 0.767 & 0.089 & 0.606 \\

HNB-linear 
& HNB (True) & linear (Wrong) 
& 0.282 & $<0.001$  & 0.002 & 0.165 
& 0.700 & $<0.001$  & 0.017 & 0.140 \\

NB-log     
& NB (Wrong) & log (Wrong)   
& 0.003 & 0.137 & 0.001 & 0.363 
& 0.044 & 0.281 & 0.001 & 0.178 \\

NB-linear  
& NB (Wrong) & linear (Wrong)
& $<0.001$  & 0.007 & $<0.001$  & 0.415 
& $<0.001$  & 0.004 & $<0.001$  & 0.368 \\
\bottomrule
\end{tabular}%
}
\end{table}

The short answer to the reviewer's question is that a functional-form
deficiency remains detectable under a misspecified response family. For
the NB-linear model, both the SW test and the AOV test against $X_1$
reject decisively ($p$-values $<0.001$ and $0.007$ for posterior
Z-residuals, respectively, with similar ISCV results), and the residual
scatterplot against $X_1$ still displays a visible and pronounced trend.

The situation is more nuanced for the NB-log model, which clearly
illustrates how a likelihood misspecification can impact functional form.
While this model uses the logarithmic form that was correct for the HNB
process, fitting a single NB family absorbs the $X_1$-dependent hurdle
probability into the NB mean structure. Consequently, the log-linear form
alone fails to correctly specify the conditional mean for the NB model.
As expected, the SW test flags the overall distributional mismatch strongly
($p$-values $0.003$ and $0.044$). The AOV test against $X_1$ yields moderate
median upper-bound $p$-values ($0.137$ and $0.281$). While these do not strictly
cross the $0.05$ threshold here, they are conservative by design: the
replicated $p$-values actually concentrate between $0.05$ and $0.10$ (and as an
aside, frequently drop below $0.05$ in supplementary simulations with other
datasets). This moderate AOV signal correctly detects the induced mean
distortion rather than acting as a false alarm, a finding corroborated by a
clear near-linear upward trend in the corresponding residual scatterplot.

More broadly, the reviewer's question highlights the intricacies of
disentangling different sources of model error, a topic we regard as a
very important avenue to pursue. Our general finding is that the SW test
is highly sensitive to response-family misspecifications that alter the
global residual distribution, whereas residual-versus-covariate
scatterplots and grouped AOV tests are tuned to localized functional-form
errors. However, these sources are not entirely separable: as demonstrated
by the NB-log example, a misspecified response family can distort the
implied mean structure, causing family errors to legitimately propagate
into covariate-specific tests. Evaluating these diagnostics jointly helps
practitioners pinpoint the nature of a model's deficiency, but isolating
the exact source remains challenging. A robust diagnostic workflow thus
requires combining various graphical and numerical residual tools,
supplemented by model comparison metrics such as information criteria. We
are also working on extending the proposed Z-residuals to compute
residuals for different sub-models (e.g., the zero-count component of
hurdle models). We have added a discussion of these points in Section~S1.2.
Considering page limits, the main text remains focused on comparing the
proposed Z-residual methods with PPC, HPC, and SW-based tests.

\medskip

\mybox{gray}{\textbf{Question 3(b)} Would it make sense to compare the $z$-residuals against a covariate that has not been used by the model?}

\paragraph{Response:}

Thank you for this helpful question. Yes, examining Z-residuals against a covariate not included in the fitted model is a valid and highly useful strategy for discovering omitted predictors, in direct analogy with classical regression diagnostics for detecting omitted predictors. A systematic trend, curvature, or grouped mean shift in the Z-residuals against an unused covariate may indicate a missing covariate, interaction, or nonlinear effect, and can be assessed with a scatterplot with a smooth trend or a grouped ANOVA test.

We have added a discussion of this use in Section 2.5 of the revised manuscript, where we also note that it should be treated as exploratory: when many predictors are screened, multiple testing bias should be taken into account.

\medskip

\mybox{gray}{\textbf{Question 4.} The general style of the writing and the way the results are presented and discussed need improvement.}

\paragraph{Response:} Thank you for the following specific suggestions which have been incorporated in the revised manuscript and have greatly strengthened this paper.

\mybox{gray}{\textbf{Question 4(a)} I believe the authors should present the exact steps of their approach compactly in a simple algorithm. Since the authors are providing a general procedure, having a clear algorithm would significantly facilitate revisiting the paper, if one needs it.}


\paragraph{Response:}
Thank you for this helpful suggestion. We agree that the exact steps of the proposed approach should be easy for readers to revisit and implement, and we have addressed this point in two ways.

First, \textbf{Table~1} in Section~2.6 already presents the algorithm: it lays out, step by step, the computing workflow of the proposed Z-residual diagnostics alongside that of PPC with parameter-wise Z-residuals, highlighting the key operational distinction between the proposed ``summarize-first, then-test'' workflow and the ``test-first, then-summarize'' workflow of PPCs. To make this algorithmic summary more visible to readers, we have revised the table caption to ``Computational Procedures of PPC and Z-residual Diagnostics'' and the section title accordingly.

Second, to facilitate implementation, we have added a new section ``Data and
Software Availability,'' which points readers to the accompanying
\texttt{Zresidual} R package. The package includes:
\begin{itemize}
\item a generic function \texttt{Zresidual} for computing Z-residuals, with
built-in support for \texttt{brms} model objects, Cox proportional hazards
and parametric survival models from the \texttt{survival} package, and
standard GLMs from \texttt{stats};
\item high extensibility to custom model-fitting objects via user-supplied
RSP functions, through the \texttt{log\_cdf} and \texttt{log\_pmf}
arguments;
\item a comprehensive suite of graphical and numerical diagnostic tools for
evaluating Bayesian regression models;
\item the processed dolphin-sighting data and 
\item compiled HTML vignettes and Quarto sources
for reproducing the results of all the simulation studies and the
real-data analysis as reported in this paper, as well as many other vignettes such as an animated illustration of the parameter-wise Z-residual.
\end{itemize}
An anonymized version is available at
\url{https://figshare.com/s/00ac6b406e451c76e8e1} for the double-blind
review; the package will be published publicly on GitHub and submitted to
CRAN upon acceptance. Once this package is downloaded, one can start with the file "index.html" in ``docs'' directory to preview the pre-compiled HTML vignettes. 

\medskip

\mybox{gray}{\textbf{Question 4(b)} The paper contains a lot of abbreviations and combinations of hyphens (PPP, PPC, NB, HNB, Post-MSP, $\mathrm{RSP}_i^{CV}$, $Z^{\mathrm{ISCV-RSP-SW}}$, etc.), which are sometimes hard to follow. While these are summarized on the title page and have a rationale, the readability can be improved by limiting their use.}

\paragraph{Response:}
We agree that the original manuscript used too many abbreviations and that this affected readability. In the revised manuscript, we have reviewed all abbreviations and reduced their use where possible. In particular, low-frequency abbreviations, such as CDF, PMF, CPO, CV, LOOCV, and PSIS-CV, are now written out in full or avoided when they are not needed. We also write terms such as goodness-of-fit, posterior predictive p-value, probability integral transform, hurdle Poisson, and Bartlett test in full in the main text when space allows.

At the same time, some abbreviations are difficult to remove completely because they refer to core diagnostic methods or model classes that are used repeatedly in formulas, figures, and tables. For example, RSP, PPC, ISCV, NB, HNB, SW, and AOV appear frequently in method labels and figure/table annotations. Writing these terms in full throughout the manuscript would substantially increase the length of the paper and make several figure labels and tables difficult to read. Therefore, we retain a small set of core abbreviations, define them clearly at first use, and use them consistently throughout the manuscript. We have also revised figure captions, table captions, and section headings to make them more self-contained, so that readers do not need to repeatedly refer back to the abbreviation list.

We have revised the notation and wording throughout the main text, especially Sections 2--4, and updated figure and table captions to reduce low-frequency abbreviations and improve readability.

\medskip

\mybox{gray}{\textbf{Question 4(c)} The paper only provides a high-level summary of the theoretical results. It would be better to report the core of some of the statements (such as Theorem S2.2 and S2.5) in the main manuscript, while leaving the list of assumptions in the supplement.}

\paragraph{Response:}
Thank you for this helpful suggestion. Following it, we now restate the two fundamental theorems in Section~2.7 of the revised manuscript: the theorem establishing the total variation convergence of the LOO predictive distribution to the true data-generating distribution, and the theorem establishing the asymptotic independence and uniformity of the cross-validated RSPs (Theorems~S2.2 and S2.5 in the original submission; their numbers have changed in the revision because additional sections were added to the Supplementary Materials). The theorems are stated in the main text with informal one-line summaries of the regularity conditions, while the full lists of assumptions and the proofs remain in the supplement, and the main-text numbering is kept consistent with the supplement. We agree that these technical statements substantially clarify the verbal description of the asymptotic results and make the theoretical foundation of the proposed diagnostics easier to grasp without consulting the supplement.

\setlength{\bibsep}{0.2em}
\bibliographystyle{agsm}
\bibliography{ref}

\end{document}